\graphicspath{{chapters/chapter2/img}}
\chapter{Capitolo 2} \label{chapter2}
Lo stato dell'arte relativo alla sincronizzazione e alla ricostruzione di scene dinamiche multi-camera ha subito una trasformazione radicale grazie all'avvento dei modelli di fondazione visiva e di nuove tecniche di ottimizzazione geometrica. La sincronizzazione video è stata esplorata da più prospettive. I metodi basati sulla geometria, come quelli di Albl et al. e Li et al., stimano offset temporali usando la geometria epipolare, ma tipicamente assumono scene statiche o punti di vista fissi. Gli approcci human-centric usano la pose umana come segnale di sincronizzazione, ma sono limitati dall'accuratezza della stima della pose, dal numero di persone nella scena, e faticano in scenari senza attività umana prominente. Gli approcci basati sull'audio usano indizi audio per la sincronizzazione, ma funzionano solo in ambienti silenziosi e non sono
generalmente applicabili a contesti reali in cui l'audio è rumoroso o non disponibile. I metodi basati
sull'apprendimento come Sync-NeRF, ottimizzano congiuntamente pose e offset temporali, ma
spesso sono vincolati ad ambienti o tipi di oggetti specifici. Per la Multi-View Structure-from-Motion ci sono tecniche Structure-from-Motion (SfM), come COLMAP, che hanno significativamente avanzato le pipeline di ricostruzione 3D producendo pose accurate delle fotocamere da immagini e video multi-vista. Modelli più recenti come HLOC, Hierarchical-Localization, e VGGT, Visual Geometry Grounded Transformer, si basano su questi progressi usando feature basate sull'apprendimento e transformer per gestire il matching su larga scala e la stima della pose. Sebbene questi metodi raggiungano forti prestazioni nella stima della geometria delle fotocamere, sono carenti nella sincronizzazione di scene dinamiche, poiché si basano principalmente su indizi visivi statici. Per il tracking e le corrispondenze ci sono modelli come CoTracker tracciano punti densamente nel tempo, offrendo una forte coerenza temporale. Tuttavia, non modellano le corrispondenze spaziali tra punti di vista diversi. D'altro canto, MASt3R si concentra sul matching spaziale e la ricostruzione stereo, fornendo corrispondenze cross-view dense, ma non gestisce la dinamica temporale, specialmente in scene in movimento. La segmentazione si è spostata dai modelli addestrati su classi specifiche come persone o veicoli, a sistemi universali. Il modello SAM 2, Segment Anything Model 2, è l'attuale punto di riferimento ed è un modello fondamentale per la risoluzione del problema della segmentazione visiva guidata da prompt in immagini e video tramite una memoria di streaming che permette di mantenere la coerenza temporale anche in sequenze lunghe e complesse. Framework come DEVA hanno ulteriormente raffinato questo approccio, introducendo un sistema di propagazione bidirezionale che de-normalizza gli errori dei modelli a frame singolo.